\section{The degree-\(21\) divisor and terminal certificates}
\label{app:degree-twenty-one}

The body develops the lattice-gap quotient covers.  Here we record the
complementary calculation on the original lower face and the exact
certificates obtained after substituting its five dessins into the two
surviving Newton supports.

\subsection{The lower face}

For either surviving support in the Guccione--Guccione--Horruitiner--Valqui
reduction \cite{guccioneEtAl2022degree108}, the transformed bracket is
\[
[P,Q]=X^2.
\]
For the primitive toric valuation
\[
\nu(X)=-2,\qquad \nu(Y)=1,
\]
put \(z=XY^2\).  The lower faces have the form
\[
P_{\mathrm{face}}=Xp(z),\qquad
Q_{\mathrm{face}}=X^2Yq(z),
\]
where \(\deg p=7\), \(\deg q=10\), and
\begin{equation}
\label{eq:degree21-face}
pq+2zpq'-3zp'q=1.
\end{equation}

\begin{proposition}[The degree-\(21\) Belyi map]
\label{prop:degree21-belyi}
The rational function
\[
\tau(z)=z\frac{q(z)^2}{p(z)^3}
\]
is a degree-\(21\) Belyi map with passport
\[
(2^{10}1),\qquad(3^7),\qquad(17\,1^4).
\]
There are exactly five connected dessin isomorphism classes with this
passport.  They form one arithmetic orbit over an irreducible quintic field
and have monodromy \(A_{21}\).
\end{proposition}

\begin{proof}
Equation \eqref{eq:degree21-face} implies that \(p,q\) are coprime and
nonzero at zero.  Differentiation gives
\[
\tau'(z)=\frac{q(z)}{p(z)^4}.
\]
Thus the ten roots of \(q\) give double zeros, the seven roots of \(p\)
give triple poles, and infinity has ramification index \(17\).  The total
ramification is
\[
10+14+16=40=2\cdot21-2,
\]
so there is no additional branch value.

The exact Murnaghan--Nakayama character calculation in \(S_{21}\), followed
by the class-size factor, gives \(5\cdot21!\) labeled triples.  Each is
transitive, and the deck group is trivial, so there are five isomorphism
classes.  Exact coefficient reconstruction places all five over one
quintic field; exact permutation certificates give monodromy \(A_{21}\).
The enumeration and reconstruction are computer-assisted.
\end{proof}

\subsection{Identification with Borisov's divisor}

\begin{theorem}[Valuation over \(F_{-5}\)]
\label{thm:borisov-fminusfive}
Let \(E_\nu\) be the source divisor defined by \(\nu\).  Its target center is
Borisov's divisor \(F_{-5}\).  The normal ramification index and residue
degree are
\[
(e,f)=(1,21).
\]
Consequently neither surviving Newton support can realize Borisov's
Three-dessin framework.
\end{theorem}

\begin{proof}
Set \(t=Y\), so \(x=t^2/z\) in the original source coordinates.  The source
volume form is
\[
dx\wedge dy=t^{-6}z^2\,dt\wedge dz.
\]
After adding the reduced boundary coefficient, the augmented-canonical label
is \(-5\).

The face expansions are
\[
P=t^{-2}zp(z)+\text{higher \(\nu\)-terms},\qquad
Q=t^{-3}z^2q(z)+\text{higher \(\nu\)-terms}.
\]
On the target put
\[
s=\frac P Q,\qquad \lambda=\frac{Q^2}{P^3}.
\]
Then \(s\) is a uniformizer and
\[
dP\wedge dQ=-s^{-6}\lambda^{-3}\,ds\wedge d\lambda,
\]
so the target label is also \(-5\).  Pullback gives
\[
\frac P Q
=t\left(\frac{p(z)}{zq(z)}+O(t)\right),
\]
hence \(e=1\).  The residue of \(\lambda\) is
\[
\operatorname{res}_{E_\nu}(\lambda)
=z\frac{q(z)^2}{p(z)^3}=\tau(z),
\]
so \(f=21\) by \cref{prop:degree21-belyi}.

In Borisov's target coordinates \((y_1,y_2)=(Q,P)\), the divisor
\(F_{-5}\) is characterized by valuations \((-3,-2)\) and by the
nonconstant residual parameter \(y_1^2/y_2^3\).  Those are exactly the
values above.  Later infinitely-near divisors with the same two numerical
orders have constant residue of this parameter, so they are distinguished.

Borisov's Three-dessin framework gives residue degree \(16\) on its unique
source divisor above \(F_{-5}\) \cite{borisov2020frameworks}.  The forced
divisor has residue degree
\(21\), a contradiction.  This excludes the named framework, not the two
supports themselves.
\end{proof}

\begin{corollary}
Any plane counterexample of maximum coordinate degree below \(125\) has
generic degree at least \(21\).
\end{corollary}

\begin{proof}
The published below-\(125\) reduction leaves the two supports above, and the
generic-degree formula contains the contribution \(ef=21\) from
\cref{thm:borisov-fminusfive}.
\end{proof}

\begin{question}
Is \(E_\nu\) the only source divisor over the generic point of \(F_{-5}\)?
The calculation \(f=21\) is one residue-degree contribution; cancellation
valuations or infinitely-near divisors could contribute additional sheets.
\end{question}

\subsection{Exact support certificates}

Substituting the five reconstructed lower faces into the truncated Newton
support gives seven compatibility equations over the quintic field.  A
nonzero exact resultant, certified by good reduction, eliminates every
dessin.  Thus the truncated support is empty in characteristic zero.

For the full support, exact layer elimination and normalization reduce the
problem to fifteen equations in five variables over the same quintic field.
Six selected equations already generate the unit ideal.

\begin{theorem}[Terminal unit-ideal certificate]
\label{thm:terminal-unit}
For the stored normalized full-support system, the six selected obstruction
equations generate the unit ideal over the quintic coefficient field.
\end{theorem}

\begin{proof}[Exact computer-assisted certificate]
The Macaulay multiplication map has \(10{,}824\) columns.  The certificate
specifies \(7{,}121\) rows and \(7{,}121\) pivot columns.  Reducing this
exact minor at a good prime and at the selected quintic embedding gives
\[
\det=859\pmod{2053}.
\]
The determinant is therefore nonzero in characteristic zero.  The
certificate package regenerates all fifteen exact equations, rebuilds the
matrix, and replays the fixed pivot minor.
\end{proof}

Combining the truncated resultant and
\cref{thm:terminal-unit} eliminates both encoded support alternatives after
the degree-\(21\) face substitution.  This statement is exact for the stored
systems.  Turning it into a stand-alone global below-\(125\) proof requires
a line-by-line audit that the published Newton reduction, the five faces,
all saturations and normalizations, and the layer systems exhaust every
candidate.

There is also a smaller geometric compression.  Five partial-gluing
equations have mixed volume \(296\); on the certified
\(296\)-dimensional modular finite algebra, the sixth obstruction is a unit,
and its rational norm over the five split embeddings is
\[
51\pmod{2053}.
\]
This is already a compact representative of the large certificate.  To use
it as an independent characteristic-zero proof, one still needs a
publication-level toric-saturation lemma showing that the finite \'{e}tale
special fiber exhausts the complete toric intersection.

\subsection{Credit for the degree-\(125\) conclusion}

MathOverflow user \texttt{ratto3423} publicly announced the conclusion that
every characteristic-zero plane counterexample has maximum coordinate degree
at least \(125\), based on a computer-assisted elimination of the same two
supports \cite{ratto3423degree125}.  We credit that public theorem
announcement to \texttt{ratto3423}.  The public answer does not disclose
enough method to compare its computation with the certificates above.

Accordingly, our claim is not priority for the degree-\(125\) conclusion.
The potential contribution of this project is the explicit degree-\(21\)
boundary description and an independently auditable terminal certificate
architecture.  The current certificate is conditional only at the upstream
exhaustiveness layer described above, not at its final linear-algebra step.
